\chapter{函数空间 \texorpdfstring{$C(K)$}{C(K)}}
\label{chap:continuous-function-space-solutions}

\begin{solution}{ck-sup-norm-details}
紧致性和连续性保证 $|f|$ 取得有限最大值，所以范数定义良好。若
$\norm f_\infty=0$，则每点 $|f(x)|=0$，故 $f=0$；零函数的上确界范数
也等于零。齐次性由
$\max|\lambda f|=|\lambda|\max|f|$ 得到。逐点三角不等式给出
$|f(x)+g(x)|\le\norm f_\infty+\norm g_\infty$，再取最大值得三角不等式。
若 $K$ 非紧，连续函数可能无界，例如 $f(x)=x$ 在 $\R$ 上，所以上确界未必
有限；也可能只有上确界而不取得最大值。
\end{solution}

\begin{solution}{ck-uniform-cauchy-criterion}
若 $f_n\to f$ 于上确界范数，则
$\norm{f_n-f_m}_\infty\le\norm{f_n-f}_\infty+\norm{f_m-f}_\infty$，故为
Cauchy 列。反之，定理 \ref{thm:ck-banach} 说明 $C(K)$ 完备，所以每个
上确界范数 Cauchy 列收敛于某个 $f\in C(K)$。由命题
\ref{prop:sup-norm-on-ck}，这种收敛就是一致收敛。
\end{solution}

\begin{solution}{integral-norm-triangle-minkowski-finite}
$p=1$ 时由 $|f+g|\le|f|+|g|$ 积分即得。$p=2$ 时先由
有限和的 Cauchy--Schwarz 不等式，令同一分割上的 Riemann 和趋于积分，得
$\int_a^b|f g|\le\norm f_2\norm g_2$。因此实值、复值情形均有
$|\int_a^b f\bar g|\le\norm f_2\norm g_2$。再展开
\[
 \norm{f+g}_2^2=\norm f_2^2+2\operatorname{Re}\int f\bar g+\norm g_2^2
 \le(\norm f_2+\norm g_2)^2.
\]
两边取非负平方根即得三角不等式。
\end{solution}

\begin{solution}{spike-integral-norms}
在中心两侧以 $t$ 表示到端点的线性坐标，一边的积分为
$\int_0^{1/n}(nt)^pdt=1/(n(p+1))$。两边相加得
\[
 \norm{f_n}_p=\left(\frac{2}{n(p+1)}\right)^{1/p}\to0.
\]
对 $x\ne1/2$，充分大 $n$ 时 $x$ 落在支集外，故 $f_n(x)\to0$；但
$f_n(1/2)=1$，所以在中心点不收敛于零。
\end{solution}

\begin{solution}{c-not-complete-lp}
沿用正文 $n\ge2$ 时的 $f_n$ 与阶跃函数 $h$。二者只在长度 $2/n$ 的区间上不同且差的
绝对值不超过 $1$，故
$\int_0^1|f_n-h|^p\le2/n\to0$。正文的有限和证明同样适用于分段连续
函数，故可对 $f_n-h$ 使用积分三角不等式，得 $(f_n)$ 是 $p$-范数
Cauchy 列。若它还在 $C([0,1])$ 中收敛于连续 $g$，再次用三角不等式
得 $\int|g-h|^p=0$。若左半区某点 $g\ne0$，连续性给出一段正长度区间
上的正下界，导致积分为正；右半区对 $g-1$ 同理。因此 $g$ 左半区恒为零、
右半区恒为一，在 $1/2$ 处不能连续。所有有限 $p$ 都不完备。
\end{solution}

\begin{solution}{evaluation-functional-norm}
上界为
$|T(f)|\le\sum_j|a_j|\norm f_\infty$。下面构造达到这个上界的函数。
紧致 Hausdorff 空间正规。由于点 $x_1,\ldots,x_m$ 两两不同，可以先取
两两不交的开邻域 $U_j$，再取开集 $V_j$ 使
$x_j\in V_j\subseteq\overline{V_j}\subseteq U_j$。对闭集
$\overline{V_j}$ 与 $K\setminus U_j$ 使用 Urysohn 引理，得到
$\varphi_j\in C(K;[0,1])$，使
\[
 \varphi_j|_{\overline{V_j}}=1,\qquad
 \varphi_j|_{K\setminus U_j}=0.
\]
令 $c_j=\overline{a_j}/|a_j|$（实数情形取符号；$a_j=0$ 时令 $c_j=0$），并置
$f=\sum_jc_j\varphi_j$。每点至多有一个 $\varphi_j$ 非零，所以
$\norm f_\infty\le1$，且
$f(x_j)=c_j$。于是
$T(f)=\sum_j|a_j|$，故
$\norm T=\sum_j|a_j|$。
\end{solution}

\begin{solution}{weighted-integral-functional-norm}
上界已由 $|I_g(f)|\le\norm f_\infty\int|g|$ 得到。取
$f_\eta=\overline g/\sqrt{|g|^2+\eta^2}$，则 $f_\eta$ 连续且范数不超过
$1$，并且
$I_g(f_\eta)=\int |g|^2/\sqrt{|g|^2+\eta^2}$。对 $r\ge0$，
\[
 0\le r-\frac{r^2}{\sqrt{r^2+\eta^2}}\le\eta,
\]
故两积分之差不超过 $(b-a)\eta$，从而前者趋于 $\int|g|$。范数至少为
该值。实值例子 $g(x)=x$（$-1\le x\le1$）的符号函数在零点跳跃，
无论如何补定义该点的值都不连续，不能直接作为 $C([-1,1])$ 中的测试函数。
有零点本身并不必然造成这种困难，例如 $g(x)=x^2$ 可以取常函数 $f=1$
达到范数。复值情形对应的是共轭相位 $\bar g/|g|$，同样须检验其能否
在零点处连续延拓。
\end{solution}

\begin{solution}{point-evaluation-l1-discontinuous}
固定 $x_0\in(0,1)$。对 $n\ge1$，置
$r_n=\min\{x_0,1-x_0\}/(n+1)$，取
$f_n(x)=\max\{1-|x-x_0|/r_n,0\}$。三角形完整落在区间内，故
$f_n(x_0)=1$，而 $\norm{f_n}_1=r_n\to0$。若点值泛函连续，
$f_n\to0$ 应推出 $f_n(x_0)\to0$，矛盾。对 $x_0=0$ 取
$f_n(x)=\max\{1-nx,0\}$，其积分为 $1/(2n)$；对 $x_0=1$ 作反射即可。
\end{solution}

\begin{solution}{ck-lattice-lipschitz}
反三角不等式给出 $||f(x)|-|g(x)||\le|f(x)-g(x)|$，取上确界即得
$\norm{|f|-|g|}_\infty\le\norm{f-g}_\infty$。利用
$f\vee g=(f+g+|f-g|)/2$ 可得
\[
 \norm{(f\vee g)-(f'\vee g')}_\infty
 \le\norm{f-f'}_\infty+\norm{g-g'}_\infty,
\]
最小值同理。因此两种二元格运算连续。
\end{solution}

\begin{solution}{neumann-inverse-function-algebra}
记 $r=\norm f_\infty<1$。由次乘性，对 $M>N\ge0$ 有
\[
 \norm{\sum_{n=N+1}^M f^n}_\infty
 \le\sum_{n=N+1}^M r^n\le\frac{r^{N+1}}{1-r}\longrightarrow0.
\]
所以部分和是 $C(K)$ 中的 Cauchy 列，完备性保证它收敛于 $S\in C(K)$。
部分和 $S_N=1+f+\cdots+f^N$ 满足
$(1-f)S_N=1-f^{N+1}$。令 $N\to\infty$，右端一致趋于 $1$，连续乘法给出
$(1-f)S=1$；交换性也给出 $S(1-f)=1$。故 $S=(1-f)^{-1}$。
\end{solution}

\begin{solution}{ck-infinite-dimensional}
从无限 $K$ 中任取互异点 $x_1,\ldots,x_n$。有限点集闭，利用正规性可构造
$f_j\in C(K)$，使 $f_j(x_j)=1$ 而 $f_j(x_i)=0$（$i\ne j$）。若
$\sum c_jf_j=0$，在 $x_i$ 取值得 $c_i=0$，故这些函数线性无关。$n$ 任意，
所以 $C(K)$ 不可能有限维。
\end{solution}

\begin{solution}{disjoint-tent-no-uniform-subsequence}
对 $n\ge1$，取
$I_n=[(3/4)2^{-n},(5/4)2^{-n}]$。后一段的右端点
$(5/8)2^{-n}$ 小于前一段的左端点，故这些区间两两不交且都在 $(0,1)$ 中。
令 $f_n(x)=\max\{1-4\cdot2^n|x-2^{-n}|,0\}$。它是支集为 $I_n$、
高为 $1$ 的三角形函数。则
$\norm{f_n}_\infty=1$；若 $m\ne n$，在 $f_n$ 的峰点有 $f_m=0$，故
$\norm{f_n-f_m}_\infty=1$。因此该序列没有 Cauchy 子列，也没有一致收敛
子列。
\end{solution}

\begin{solution}{finite-dimensional-unit-ball-compact}
零维空间的闭单位球为单点。设 $n\ge1$，在 $X$ 选基 $e_1,\ldots,e_n$，令
$T:\mathbb F^n\to X$，$T(a)=\sum a_ie_i$。三角不等式给出
\[
 \norm{T(a)-T(b)}\le\sum_i|a_i-b_i|\norm{e_i}
 \le\Bigl(\sum_i\norm{e_i}^2\Bigr)^{1/2}\norm{a-b}_2.
\]
故 $T$ 连续。复坐标空间视为实 $2n$ 维空间时同一估计成立。
单位坐标球面紧，连续函数
$a\mapsto\norm{T(a)}$ 在其上取得正的最小值 $m$ 和有限最大值 $M$；最小值
为正是因为基线性无关。故
$m\norm a_2\le\norm{T(a)}\le M\norm a_2$。$X$ 的闭单位球在 $T^{-1}$ 下
对应欧氏空间中的闭有界集，因而紧；连续像仍紧。这同时证明了所需的范数
等价，无须预先假定。
\end{solution}

\begin{solution}{riesz-lemma-closedness}
在 $X=C([0,1])$ 中，令 $M$ 为所有分段线性函数组成的子空间。它是真
子空间；另一方面，连续函数一致连续，在足够细的分割节点上作线性插值，
便可在上确界范数下任意逼近原函数。因此 $M$ 在 $X$ 中稠密。任意
$x\in X$ 到 $M$ 的距离都是零，不可能大于给定 $\alpha>0$。Riesz 引理
证明中 $d(z,M)>0$ 正是由 $M$ 闭且 $z\notin M$ 得到；删去闭性后这一步
失败。
\end{solution}

\begin{solution}{compact-operator-first-example}
有界性由
\[
 |Tf(x)|\le\norm f_\infty\int_0^1|k(x,y)|dy
 \le\norm k_\infty\norm f_\infty
\]
得到，故 $\norm T\le\norm k_\infty$。若 $\norm f_\infty\le1$，则
$|Tf(x)|\le\norm k_\infty$；并且
\[
 |Tf(x)-Tf(x')|\le\int_0^1|k(x,y)-k(x',y)|dy.
\]
$k$ 在紧方形上一致连续，所以右端对整个单位球一致趋于零。像集因而一致
有界且等度连续；下一章的 Arzelà--Ascoli 定理将说明其闭包紧致。
\end{solution}
